\chapter{Seasonality and Calendar Effects}\label{ch:s1:seasonality-and-calendar-effects}

Hundreds of \omterm{def:s1:seasonality-and-calendar-effects:calendar}{calendar effects} have been published: Mondays, Januaries, the turn of the month, the days before holidays, the halves of the month, the
seasons. Search enough of them and some will look significant in any sample. Lakonishok and Smidt, with ninety years of the Dow Jones Industrial
Average, found returns that were persistently anomalous around the turn of the week, the month, the year and holidays; natural gas is dearer in winter
for a reason everybody knows. On this chapter's synthetic market two \omterm{def:s1:seasonality-and-calendar-effects:calendar}{calendar effects} are planted among a hundred rules: the turn of the month survives
every correction for multiple testing, the pre-holiday effect is too small to be found in thirty years, and the uncorrected search reports eight
effects that do not exist. The build is \texttt{firm.seasonal}.

\begin{definition}[Calendar effect, return seasonality]\label{def:s1:seasonality-and-calendar-effects:calendar}
A \emph{calendar effect}\index{calendar effect} is a difference between average returns on days, weeks or months defined by the calendar alone (a
weekday, a day of the month, a month, the days around a holiday) and average returns on other days. \emph{Return seasonality}\index{return
seasonality} is the recurrence of a security's relative return at the same point of each year, such as the same calendar month.
\end{definition}

\section{Commodity seasonality with a physical reason}

Some seasons are physical. Natural gas is burned for heating in winter and stored in summer. From the Henry Hub spot price, 2000 to 2026 (EIA, via
FRED), the log price's deviation from its centred twelve-month average is $+8.8\%$ in January and $+7.3\%$ in December, and $-5.3\%$ in March and
$-4.2\%$ in April, as storage refills (\cref{fig:s1:seasonality-and-calendar-effects:months}). That pattern is not a trade. Futures for January delivery
are priced for January: the curve anticipates the season, as \texttt{firm.synthfut}'s seasonal commodities do by construction, so a futures position
earns the season only if it was mispriced. What a trader can hold is the risk: a winter--summer spread that pays when winter turns out colder, or
storage that turns out tighter, than the curve priced.

Crude oil is the counterexample: its demand is less seasonal and storage is ample. Testing each calendar month's average rolled WTI return against the
other months', 1986 to 2024, the lowest are October ($-2.8\%$ a month, $t = -1.71$) and November ($-4.2\%$, $t = -2.17$, a raw p-value of 0.03). Twelve
months were tested; after Bonferroni's or Benjamini and Hochberg's correction the smallest adjusted p-value is 0.36. There is nothing there.

\begin{omfigure}
\begin{tikzpicture}
\begin{axis}[width=10.5cm,height=5cm,ybar,bar width=8pt,xlabel={\small calendar month},ylabel={\small deviation from trend (\%)},
  tick label style={font=\footnotesize},xtick={1,2,3,4,5,6,7,8,9,10,11,12},xmin=0.4,xmax=12.6,ymin=-7,ymax=11,grid=major,grid style={gray!25},
  table/col sep=comma]
\addplot[fill=omThm!60,draw=omThm] table[x=month,y=gas_pct]{figdata/strategies-1/23-seasonality-and-calendar-effects/months.csv};
\end{axis}
\end{tikzpicture}

\omcaption{The Henry Hub natural gas spot price's seasonal profile: the average deviation of the monthly log price from its centred twelve-month
average, by calendar month, January 2000 to July 2026 (EIA via FRED, series MHHNGSP). Data: \texttt{s1\_seasonal.gas\_profile}.}\label{fig:s1:seasonality-and-calendar-effects:months}
\end{omfigure}

\section{Turn-of-the-month and holiday effects}

\begin{definition}[Turn-of-the-month effect]\label{def:s1:seasonality-and-calendar-effects:tom}
The \emph{turn-of-the-month effect}\index{turn-of-the-month effect} is the tendency of stock market returns to be higher on the last trading day of a
month and the first few days of the next than on other days; month-end flows (salaries, pension contributions and fund rebalancing invested at the
start of the month) are the candidate mechanism.
\end{definition}

The chapter's test is built to be honest about searching. \texttt{firm.seasonal} lays out a stylised calendar (twelve months of 21 trading days, weeks
of five days, nine holidays a year) under thirty years of \texttt{firm.synthfut}'s equity class (the mean of its ten index futures, 14.7\% volatility a
year). Two effects are planted: 8 basis points a day on the last day and first three days of each month, and 15 basis points on each day before a
holiday. A hundred rules are generated (\cref{lst:s1:seasonality-and-calendar-effects:rules}): five weekdays, twelve months, 21 days of the month, twelve
turn-of-the-month windows, the days before and after holidays, weeks of the month, fortnights of the year, first and last days of the month, and a few
seasonal favourites (``May to October'', ``December and January''). Each rule compares the mean return on its days with the mean on the others. Of the
hundred, 23 overlap the turn-of-the-month days enough to carry that effect (an expected difference of at least 3 basis points), 2 carry the
pre-holiday effect, and 75 carry nothing.

\begin{center}
\small
\begin{tabular}{@{}lccc@{}}
\toprule
correction (5\%) & turn-of-month rules found & pre-holiday rules found & rules with no effect found\\
\midrule
none & 19 of 23 & 0 of 2 & 8 of 75\\
Bonferroni & 6 & 0 & 0\\
Holm & 6 & 0 & 0\\
Benjamini--Hochberg & 11 & 0 & 3\\
Benjamini--Yekutieli & 6 & 0 & 0\\
\bottomrule
\end{tabular}
\end{center}

The turn of the month, planted on 48 days a year, is found by every correction: its best rule, the last day and first three days, has $t = 4.4$
(\cref{fig:s1:seasonality-and-calendar-effects:rules}). The pre-holiday effect, on nine days a year, has $t = 1.5$: real, larger per day, and invisible,
because nine days a year for thirty years is 270 observations. Without correction the search reports eight false effects, among them two weekdays, three
days of the month, a week of the month and two fortnights, each of which would make a paper's table; Benjamini--Hochberg, which controls the share of false discoveries
rather than the chance of any, keeps three of them; the family-wise corrections keep none. Book~4, chapter~12 derived the corrections; here they decide
what a researcher believes.

\begin{omfigure}
\begin{tikzpicture}
\begin{axis}[width=11cm,height=5.4cm,xlabel={\small rule, ranked by $|t|$},ylabel={\small $|t|$},tick label style={font=\footnotesize},
  xmin=0,xmax=101,ymin=0,ymax=5,grid=major,grid style={gray!25},legend columns=3,legend style={font=\scriptsize,at={(0.5,-0.3)},anchor=north,
  draw=none,fill=none,/tikz/every even column/.append style={column sep=8pt}},unbounded coords=discard,table/col sep=comma]
\addplot[only marks,mark=*,mark size=1.6pt,omThm] table[x=rank,y=tom]{figdata/strategies-1/23-seasonality-and-calendar-effects/rules.csv};
\addplot[only marks,mark=square*,mark size=1.8pt,omDef] table[x=rank,y=pre]{figdata/strategies-1/23-seasonality-and-calendar-effects/rules.csv};
\addplot[only marks,mark=o,mark size=1.4pt,black!60] table[x=rank,y=none]{figdata/strategies-1/23-seasonality-and-calendar-effects/rules.csv};
\addplot[black!50,dashed,forget plot] coordinates {(0,1.96) (101,1.96)};
\addplot[black,dashed,forget plot] coordinates {(0,3.48) (101,3.48)};
\legend{turn-of-month rules,pre-holiday rules,rules with no effect}
\end{axis}
\end{tikzpicture}

\omcaption{A hundred calendar rules tested on thirty years of the synthetic equity class with two planted effects: the absolute t-statistic of each rule,
ranked, by what it carries. The lower dashed line is the uncorrected 5\% threshold (1.96), the upper one Bonferroni's for a hundred tests (3.48). Data:
\texttt{s1\_seasonal.calendar\_test}.}\label{fig:s1:seasonality-and-calendar-effects:rules}
\end{omfigure}

\section{Same-month seasonality in stock returns}

Heston and Sadka found a pattern in the cross-section: stocks with high returns in a calendar month tend to have high returns in the same calendar month
in later years, at annual lags of up to twenty years, independent of size, industry, earnings announcements, dividends and fiscal year. The signal is
simple: a stock's average return in the same calendar month of past years. The chapter plants it in \texttt{firm.synthmkt}: each listing receives a
return for each calendar month, the same every year, with a cross-sectional standard deviation of 1\% a month, spread over the month's days.

The signal, averaged over the previous five years, has a monthly rank IC of 0.031 with the month's return ($t = 3.8$ over 108 months); a decile
long--short earns 0.50\% a month at a Sharpe ratio of 0.72, before costs. The planted effect is small against a monthly specific volatility of several
percent, and the signal averages only five past observations of it. A true, stable seasonal effect is a weak signal because each stock shows it once a year.

\section{What survives out of sample}

Sullivan, Timmermann and White titled their study of \omterm{def:s1:seasonality-and-calendar-effects:calendar}{calendar effects} ``dangers of data mining'', and the chapter's experiment shows the danger in
miniature: eight false discoveries out of a hundred rules, three even under a false-discovery-rate control. What survives has three things. A mechanism
that says why the days differ (month-end flows, storage, heating), stated before the test. A test that counts every rule tried, including those that
were looked at and dropped. And enough observations: an effect on nine days a year needs about half a century to reach a t-statistic of 3.48 at the
chapter's size. Most published \omterm{def:s1:seasonality-and-calendar-effects:calendar}{calendar effects} fail at least one of the three.

\section{Strategy files}

\begin{strategyfile}{Natural-gas seasonal spread}\label{strat:s1:seasonality-and-calendar-effects:gas}
\sfield{Who pays you, and why} Hedgers of winter demand who pay a premium for winter delivery beyond the expected seasonal price.
\sfield{Instruments and venues} Natural-gas futures for winter and summer delivery.
\sfield{Signal} The winter--summer spread against storage levels and its own history.
\sfield{Sizing and execution} Spread positions, small; rolled with the strip.
\sfield{Costs} Moderate; wide spreads in far months.
\sfield{How it dies} Cold winters: the risk it is paid for.
\sfield{Horizon, capacity, infrastructure} Months; storage and weather data.
\sfield{Backtest honestly} Futures prices, not the spot's seasonal profile; storage data as released.
\sfield{Sources} No performance figure verified; the spot profile from EIA data.
\end{strategyfile}

\begin{strategyfile}{Turn-of-the-month}\label{strat:s1:seasonality-and-calendar-effects:tom}
\sfield{Who pays you, and why} Month-end and month-start flows that buy equities on predictable days.
\sfield{Instruments and venues} Equity index futures.
\sfield{Signal} The calendar: the last trading day and the first three.
\sfield{Sizing and execution} Long on those days, flat otherwise; or tilt an existing book's exposure.
\sfield{Costs} Two round trips a month.
\sfield{How it dies} Arbitrage of a published pattern; changes in payroll and fund-flow timing.
\sfield{Horizon, capacity, infrastructure} Days; large capacity in index futures.
\sfield{Backtest honestly} The exact window fixed before the test; every window tried counted.
\sfield{Sources} Lakonishok and Smidt (1988); this chapter: found by every correction when planted at 8 basis points a day.
\end{strategyfile}

\begin{strategyfile}{Same-calendar-month seasonality}\label{strat:s1:seasonality-and-calendar-effects:samemonth}
\sfield{Who pays you, and why} Recurring annual flows and events specific to each firm.
\sfield{Instruments and venues} Stocks.
\sfield{Signal} A stock's average return in the same calendar month of past years.
\sfield{Sizing and execution} Monthly decile or rank long--short, neutral to factors.
\sfield{Costs} A full rebalance every month.
\sfield{How it dies} Crowding; the costs of a monthly turnover.
\sfield{Horizon, capacity, infrastructure} A month; long return histories.
\sfield{Backtest honestly} Delisted stocks included; lags of whole years only.
\sfield{Sources} Heston and Sadka (2008); this chapter: IC 0.031, 0.50\% a month decile spread before costs.
\end{strategyfile}

\begin{strategyfile}{Pre-holiday effect}\label{strat:s1:seasonality-and-calendar-effects:holiday}
\sfield{Who pays you, and why} Unclear: mood and short covering before a closure are the usual stories.
\sfield{Instruments and venues} Equity index futures.
\sfield{Signal} The trading day before an exchange holiday.
\sfield{Sizing and execution} Long for the day.
\sfield{Costs} One round trip per holiday.
\sfield{How it dies} It may never have lived: few observations a year.
\sfield{Horizon, capacity, infrastructure} A day.
\sfield{Backtest honestly} The holiday calendar as it was each year; enough years for nine events a year to mean something.
\sfield{Sources} Lakonishok and Smidt (1988); this chapter: $t = 1.5$ in thirty years when planted at 15 basis points.
\end{strategyfile}

\section{Tutorial: a handful survive}\label{tut:s1:seasonality-and-calendar-effects:tut}

\begin{tutorial}
\textbf{Goal.} Generate a hundred calendar rules, test them on a synthetic market with two planted effects, correct four ways and count true and false
discoveries; test same-month seasonality in the cross-section; read the seasons of natural gas and crude oil. \textbf{End state:} the table and the two
figures.
\begin{tutsteps}
\item \textbf{Rules and their test}.
\omcode{code/firm/seasonal/firm_seasonal.py}{39}{70}{A hundred calendar rules and the Welch test of each.}\label{lst:s1:seasonality-and-calendar-effects:rules}
\item \textbf{The experiment}: planting, testing, correcting, and classing each discovery by what it carries.
\omcode{code/strategies-1/23-seasonality-and-calendar-effects/python/s1_seasonal.py}{42}{58}{The calendar experiment.}\label{lst:s1:seasonality-and-calendar-effects:experiment}
\item \textbf{Run} \texttt{calendar\_test()}, \texttt{same\_month\_test()}, \texttt{wti\_months()} and \texttt{gas\_profile()}, and \texttt{fig\_seasonal.py}.
\end{tutsteps}
\textbf{What to change next.} Run the experiment on fifty seeds and count how often each correction finds the pre-holiday effect; add a rule family per
year of history and watch the false discoveries grow; plant a same-month effect that fades over time.
\end{tutorial}

\section{Build: calendar rules}\label{bld:s1:seasonality-and-calendar-effects:seasonal}

\begin{build}
\bfield{Purpose} A stylised calendar, a rule generator, a rule tester, seasonal profiles and the same-month signal.
\bfield{Interface} \texttt{calendar(years, holidays)}, \texttt{rules(cal)}, \texttt{score\_rules(r, rules)}, \texttt{profile(x, period)},
\texttt{same\_month(M, lags)}.
\bfield{Rules} Every rule generated is tested and counted; the same-month signal uses whole-year lags only.
\bfield{Acceptance tests} \texttt{code/firm/seasonal/tests/}: the calendar's shape and the hundred distinct rules, a planted day found with the right
p-value, profiles and the same-month signal by hand.
\bfield{Stretch} Real exchange calendars; bootstrap (max-statistic) corrections for correlated rules; seasonal decomposition of futures curves.
\end{build}

\begin{omsources}
\item S.~L.~Heston and R.~Sadka, ``Seasonality in the cross-section of stock returns'', \emph{Journal of Financial Economics} 87(2), 2008.
\item J.~Lakonishok and S.~Smidt, ``Are seasonal anomalies real? A ninety-year perspective'', \emph{Review of Financial Studies} 1(4), 1988.
\item R.~Sullivan, A.~Timmermann and H.~White, ``Dangers of data mining: the case of calendar effects in stock returns'', \emph{Journal of Econometrics}
105(1), 2001.
\item US Energy Information Administration: Henry Hub natural gas spot price (via FRED) and NYMEX WTI crude oil futures settlements.
\end{omsources}

\section{Exercises}

\begin{exercise}[$\star$]\label{exo:s1:seasonality-and-calendar-effects:1}
A hundred rules with no effect are tested at 5\%. How many false discoveries do you expect? With 75 rules that carry nothing?
\end{exercise}

\begin{exercise}[$\star$]\label{exo:s1:seasonality-and-calendar-effects:2}
What two-sided threshold does Bonferroni's correction set for a hundred tests at 5\%?
\end{exercise}

\begin{exercise}[$\star$]\label{exo:s1:seasonality-and-calendar-effects:3}
Why is the natural-gas spot price's winter premium not a trading strategy?
\end{exercise}

\begin{exercise}[$\star\star$]\label{exo:s1:seasonality-and-calendar-effects:4}
Daily volatility is $14.7\%/\sqrt{252}$. What standard error does a difference of means have for 1,440 turn-of-month days against 6,120 others, and what
t-statistic does 8 basis points give? The same for 270 pre-holiday days against 7,290 others and 15 basis points?
\end{exercise}

\begin{exercise}[$\star\star$]\label{exo:s1:seasonality-and-calendar-effects:5}
How many years would the pre-holiday effect need for its expected t-statistic to reach 3.48?
\end{exercise}

\begin{exercise}[$\star\star$]\label{exo:s1:seasonality-and-calendar-effects:6}
Why does Benjamini--Hochberg keep false discoveries that Bonferroni rejects, and when is that acceptable?
\end{exercise}

\begin{exercise}[$\star\star\star$]\label{exo:s1:seasonality-and-calendar-effects:7}
\emph{Coding.} Change the seed of \texttt{simulate\_futures} in \texttt{calendar\_test} to 8 and rerun. Which results change, and which do you expect to
stay?
\end{exercise}

\begin{exercise}[$\star\star\star$]\label{exo:s1:seasonality-and-calendar-effects:8}
\emph{Find the flaw.} ``Mondays have been negative in our data at $t = 2.4$; we short index futures every Monday.''
\end{exercise}

\section{Problem: A Handful Survive}

\begin{problem}[{Weekend problem --- a hundred calendar rules}]\label{pb:s1:seasonality-and-calendar-effects:1}
The chapter's synthetic markets, the EIA data and the public record.

\medskip\noindent\textbf{Part I --- Seasons.}
\begin{enumerate}
\item Define a \omterm{def:s1:seasonality-and-calendar-effects:calendar}{calendar effect}, the \omterm{def:s1:seasonality-and-calendar-effects:tom}{turn-of-the-month effect} and \omterm{def:s1:seasonality-and-calendar-effects:calendar}{return seasonality}.
\item Give natural gas's seasonal profile and explain it.
\item What did the WTI test by calendar month find?
\item What did Lakonishok and Smidt find?
\end{enumerate}

\medskip\noindent\textbf{Part II --- The experiment.}
\begin{enumerate}[resume]
\item Describe the calendar, the planted effects and the hundred rules.
\item How are rules classed as carrying an effect or nothing?
\item Give the discoveries at each correction.
\item Why is the pre-holiday effect never found?
\end{enumerate}

\medskip\noindent\textbf{Part III --- The cross-section.}
\begin{enumerate}[resume]
\item What did Heston and Sadka find?
\item Describe the planted same-month effect and the signal.
\item Give its IC and the decile spread.
\item Why is a stable seasonal effect a weak signal?
\end{enumerate}

\medskip\noindent\textbf{Part IV --- The verdict.}
\begin{enumerate}[resume]
\item State the \emph{named result}: the planted effects found and the false discoveries at each correction.
\item List the three things a surviving \omterm{def:s1:seasonality-and-calendar-effects:calendar}{calendar effect} has.
\item Which correction would you use for a search of a hundred rules, and why?
\item How would you backtest a turn-of-the-month strategy honestly?
\item What does ``dangers of data mining'' refer to?
\item Which strategy file has the weakest mechanism?
\item How does this chapter relate to Book~4, chapter~12?
\item In one sentence: which \omterm{def:s1:seasonality-and-calendar-effects:calendar}{calendar effects} deserve belief?
\end{enumerate}
\end{problem}

\section{Interview questions}

\begin{interviewq}[$\star$ \iqroles{researcher}]\label{iq:s1:seasonality-and-calendar-effects:1}
You find that Fridays have been positive at $t = 2.2$. What do you do next?
\end{interviewq}

\begin{interviewq}[$\star\star$ \iqroles{researcher}]\label{iq:s1:seasonality-and-calendar-effects:2}
What is the difference between controlling the family-wise error rate and the false discovery rate?
\end{interviewq}

\begin{interviewq}[$\star\star$ \iqroles{trader}]\label{iq:s1:seasonality-and-calendar-effects:3}
Why might month-end flows move equity prices, and how would you check?
\end{interviewq}

\begin{interviewq}[$\star\star$ \iqroles{developer}]\label{iq:s1:seasonality-and-calendar-effects:4}
What must a calendar in a backtesting system get right?
\end{interviewq}

\begin{interviewq}[$\star\star$ \iqroles{researcher}]\label{iq:s1:seasonality-and-calendar-effects:5}
How is seasonality in a commodity's spot price different from seasonality in its futures returns?
\end{interviewq}

\begin{interviewq}[$\star\star\star$ \iqroles{researcher}]\label{iq:s1:seasonality-and-calendar-effects:6}
A rule's days are a fraction $f$ of $n$ days, the effect is $\delta$ per rule day and daily volatility is $\sigma$. Derive the expected t-statistic of the
difference of means and the number of days needed for it to reach $z$.
\end{interviewq}
